\documentclass[11pt]{article}
\usepackage{amsmath,amssymb,amsthm}
\usepackage{hyperref}
\usepackage{graphicx}
\usepackage{xcolor}
\usepackage{booktabs}
\usepackage{longtable}
\usepackage{array}
\usepackage{calc}
\usepackage{soul}
\usepackage{geometry}
\providecommand{\tightlist}{\setlength{\itemsep}{0pt}\setlength{\parskip}{0pt}}
\geometry{margin=1in}

% Theorem environments
\newtheorem{theorem}{Theorem}[section]
\newtheorem{lemma}[theorem]{Lemma}
\newtheorem{proposition}[theorem]{Proposition}
\newtheorem{corollary}[theorem]{Corollary}
\newtheorem{conjecture}[theorem]{Conjecture}
\theoremstyle{definition}
\newtheorem{definition}[theorem]{Definition}
\newtheorem{assumption}[theorem]{Postulate}
\theoremstyle{remark}
\newtheorem{remark}[theorem]{Remark}

% Cause-plex notation shortcuts
\newcommand{\CP}{\mathcal{C}}
\newcommand{\Cstar}{\mathcal{C}^*}
\newcommand{\loopspace}{\Omega(\mathcal{C}^*, e_*)}

\title{Complex Amplitudes from Loop-Phase Consistency in the Multiway Cause-Plex}
\author{Ian Derrington}
\date{2026-03-25}

\begin{document}
\maketitle

\begin{abstract}
We prove that in any multiway cause-plex satisfying two natural conditions — bounded iterated amplitudes and causal continuity — the amplitude weights of closed causal loops must lie in $\mathrm{U}(1)$. Combined with Aaronson's minimality theorem, this forces the Feynman path integral weight $e^{iS/\hbar}$ without postulation. The imaginary unit i is not assumed: it is the unique minimal continuous group structure consistent with a stable, looping cause-plex.
\end{abstract}

\medskip

\section{Introduction}\label{introduction}

\subsection{The Problem}\label{the-problem}

Standard quantum mechanics postulates that the transition amplitude
between states is a sum over paths weighted by \(e^{iS/\hbar}\). The
imaginary unit \(i\) is taken as given. No derivation from more
primitive structure is standard. Attempts to explain it typically fall
into one of three categories:

\begin{enumerate}
\def\labelenumi{\arabic{enumi}.}
\item
  \textbf{Analogy}: Wick rotation from Euclidean to Lorentzian metric
  introduces \(i\) as a rotation between signatures. This is suggestive
  but not a derivation --- it presupposes the path integral formalism it
  is trying to ground \cite{Gorard2020}.
\item
  \textbf{Minimality arguments}: Sorkin \cite{Sorkin1994} shows quantum
  mechanics is the unique grade-2 generalization of classical
  probability theory; complex amplitudes are forced by the minimal
  relaxation of the classical sum-rule. This is a derivation from
  probability theory, not from physical causal structure.
\item
  \textbf{Uniqueness via division algebras}: Baez \cite{Baez2012} proves
  that among real, complex, and quaternionic QM, only complex QM
  satisfies both local tomography and the standard tensor product for
  composite systems. Renou et al.~(2021) experimentally confirmed that
  real QM is excluded. These are the rigorous results used in Part 3 of
  this paper. Aaronson (2004) gave an earlier informal version of the
  same argument; the Baez/Renou results supersede it.
\end{enumerate}

None derives \(i\) from the primitive causal event structure. The
present paper attempts this.

\subsection{The Cause-Plex Framework}\label{the-cause-plex-framework}

The cause-plex framework {[}Cause-Plex and Spacetime; Cause-Plex and
Quantum Mechanics{]} takes the \textbf{causal event} as the primitive: a
state transition
\(e: \mathcal{S}_{\mathrm{in}} \to \mathcal{S}_{\mathrm{out}}\) with no
assumed physics. A locally finite partial order on such events --- the
cause-plex \(\mathcal{C} = (E, \prec)\) --- gives spacetime structure,
energy, and the entity hierarchy. The \textbf{multiway cause-plex}
\(\mathcal{C}^{\ast}\) --- the collection of all causally consistent
histories --- is the substrate for quantum mechanics under Postulate Q'
(all histories coexist).

Within this framework, every stable entity is a \textbf{recurring causal
loop}: a pattern of events that returns to its initial state. Bonds,
particles, atoms, cells --- all are persistent loops in
\(\mathcal{C}^{\ast}\) at different scales {[}Part 1.5: Causors{]}. The
loop structure is not imported; it is the fundamental architectural
feature of persistent entities in the cause-plex.

The present paper asks: what amplitude weights for closed loops are
consistent with this structure?

\subsection{Main Argument in Brief}\label{main-argument-in-brief}

A closed causal loop \(\gamma\) in \(\mathcal{C}^{\ast}\) accumulates an
amplitude weight \(w(\gamma) \in G\) where \(G\) is some group. We show:

\begin{enumerate}
\def\labelenumi{\arabic{enumi}.}
\item
  \textbf{Stability forces \(|w| = 1\)}: if \(|w(\gamma)| \neq 1\),
  iterating \(\gamma\) gives amplitudes that diverge or collapse to
  zero, inconsistent with stable persistent entities (Section 3).
\item
  \textbf{Causal continuity forces \(G\) to be a topological group}:
  loops differing by one event in the branch graph must have nearby
  amplitudes (Section 4). For observer-class cause-plexes (causally
  simply-connected by Lemma 4.3), the loop space is connected (Theorem
  4.4), forcing \(G\) to be connected.
\item
  \textbf{Minimality forces \(G = \mathrm{U}(1)\)}: among connected
  compact groups acting consistently on path amplitudes with
  non-negative probabilities, U(1) is unique (Section 5, applying
  Aaronson 2004).
\end{enumerate}

The conclusion: \(w(\gamma) = e^{i\phi(\gamma)}\) for some real-valued
phase function \(\phi\) on closed loops.

\subsection{Relationship to Prior
Work}\label{relationship-to-prior-work}

{\def\LTcaptype{none} % do not increment counter
\begin{longtable}[]{@{}
  >{\raggedright\arraybackslash}p{(\linewidth - 4\tabcolsep) * \real{0.3333}}
  >{\raggedright\arraybackslash}p{(\linewidth - 4\tabcolsep) * \real{0.3333}}
  >{\raggedright\arraybackslash}p{(\linewidth - 4\tabcolsep) * \real{0.3333}}@{}}
\toprule\noalign{}
\begin{minipage}[b]{\linewidth}\raggedright
Work
\end{minipage} & \begin{minipage}[b]{\linewidth}\raggedright
What it shows
\end{minipage} & \begin{minipage}[b]{\linewidth}\raggedright
How this paper relates
\end{minipage} \\
\midrule\noalign{}
\endhead
\bottomrule\noalign{}
\endlastfoot
Feynman \& Hibbs (1965) & Path integral with \(e^{iS/\hbar}\) & We
derive the \(i\); they postulate it \\
Sorkin \cite{Sorkin1994} & QM = unique grade-2 probability theory &
Different route; our result is more structural \\
Gorard \cite{Gorard2020} & Multiway graph + complex weights \(\to\)
Hilbert space & We derive the complex weights; they assume them \\
Baez \cite{Baez2012} & Only \(\mathbb{C}\) satisfies local tomography +
standard tensor products among normed division algebras & We apply this
as Part 3; peer-reviewed, rigorous \\
Renou et al.~(2021) & Real QM experimentally falsified & Experimental
confirmation that \(\mathbb{R}\)-QM is excluded \\
Wolfram \cite{Wolfram2020} & Branchlike direction has signature \((+)\)
& Corroborating; not used directly \\
Barcelo et al.~(2001, 2006) & A-homotopy theory for graphs &
Mathematical tool for Part 2 \\
\end{longtable}
}

\medskip

\section{Definitions and Setup}\label{definitions-and-setup}

\begin{definition}[Causal event]
\label{definition:2-1}
A causal event (\emph{state couplet}) $e$ is an ordered pair $(\mathcal{S}_{\mathrm{in}}, \mathcal{S}_{\mathrm{out}})$ where $\mathcal{S}_{\mathrm{out}}$ is determined by $\mathcal{S}_{\mathrm{in}}$. No physical content is assumed.

\end{definition}

\begin{definition}[Cause-plex]
\label{definition:2-2}
The cause-plex $\mathcal{C} = (E, \prec)$ is a locally finite strict partial order on a set of causal events $E$: irreflexive, transitive, and locally finite (the interval $\{e : e_1 \prec e \prec e_2\}$ is finite for all $e_1, e_2$).

\end{definition}

\begin{definition}[Multiway cause-plex]
\label{definition:2-3}
Given initial state $\mathcal{S}_0$, the multiway cause-plex is
$$\mathcal{C}^{\ast}(\mathcal{S}_0) = \{(E_\alpha, \prec_\alpha) : \text{locally finite strict partial order consistent with } \mathcal{S}_0\}$$
the collection of all causally consistent histories.

\end{definition}

\begin{definition}[Branch graph]
\label{definition:2-4}
The branch graph $\mathcal{B}$ is the directed graph with node set $\mathcal{C}^{\ast}$ and an edge $\mathcal{C}_\alpha \to \mathcal{C}_\beta$ when $|E_\alpha \triangle E_\beta| = 1$ and $\mathcal{C}_\alpha \prec_{\mathcal{B}} \mathcal{C}_\beta$: adjacent histories differing by exactly one causal event.

\end{definition}

\begin{definition}[Causal loop]
\label{definition:2-5}
A causal loop based at event $e_0$ is a finite sequence
$$\gamma = (e_0 \prec e_1 \prec \cdots \prec e_n = e_0')$$
where $e_0'$ is in the same local state as $e_0$ — the cause-plex has returned to an equivalent local configuration. The **length** $|\gamma| = n$.

*Remark.* The precise definition of "same local state" requires specifying a local state equivalence relation. For the purposes of this paper, we work with based loops: loops based at a distinguished event $e_{\ast}$ with a fixed local state, and require $e_n$ to have identical state to $e_0$.

\end{definition}

\begin{definition}[Loop space]
\label{definition:2-6}
The loop space $\Omega(\mathcal{C}^{*}, e_{*})$ is the set of all causal loops in $\mathcal{C}^{*}$ based at $e_{*}$.

\end{definition}

\begin{definition}[Loop concatenation]
\label{definition:2-7}
For loops $\gamma_1 = (e_0, \ldots, e_n = e_0)$ and $\gamma_2 = (e_0, \ldots, e_m = e_0)$ based at $e_0$, define
$$\gamma_1 \cdot \gamma_2 = (e_0, \ldots, e_n = e_0, \ldots, e_{n+m} = e_0)$$
the concatenated path traversing $\gamma_1$ then $\gamma_2$.

\end{definition}

\begin{definition}[Amplitude assignment]
\label{definition:2-8}
An amplitude assignment is a function $w: \Omega(\mathcal{C}^{*}, e_{*}) \to G$ for some group $G$, satisfying the *composition rule*:
$$w(\gamma_1 \cdot \gamma_2) = w(\gamma_1) \cdot w(\gamma_2)$$
for all loops $\gamma_1, \gamma_2 \in \Omega(\mathcal{C}^{\ast}, e_{\ast})$.

*Remark.* The composition rule is the minimal constraint: the amplitude of two successive loops must be the product of their individual amplitudes. This is the discrete analogue of the path-integral composition property.

\end{definition}

\begin{definition}[Probability measure]
\label{definition:2-9}
Given an amplitude assignment $w$, the probability of outcome $\mathcal{S}_f$ is:
$$P(\mathcal{S}_f) = \left|\sum_{\gamma: e_{\ast} \to \mathcal{S}_f} w(\gamma)\right|^2 \cdot \mu$$
for an appropriate normalization measure $\mu$. For this to define a probability, we require $P(\mathcal{S}_f) \geq 0$ and $\sum_f P(\mathcal{S}_f) = 1$.

\end{definition}

\begin{definition}[Auto-causal density $\rho_{\mathrm{ac}}$]
\label{definition:2-10}
\end{definition}

\medskip

\section{\texorpdfstring{Part 1: Stability Forces
\(|w| = 1\)}{Part 1: Stability Forces \textbar w\textbar{} = 1}}\label{part-1-stability-forces-w-1}

\begin{proposition}[Amplitude on the unit circle]
\label{proposition:3-1}
Let $w$ be an amplitude assignment satisfying Definition 2.8. If the cause-plex supports stable persistent loops — entities that iterate indefinitely without divergence or extinction — then $|w(\gamma)| = 1$ for all $\gamma \in \Omega(\mathcal{C}^{*}, e_{*})$.
\end{proposition}

\begin{proof}
$$w(\gamma^n) = w(\gamma)^n$$
\end{proof}

Consider the behavior as \(n \to \infty\):

\begin{itemize}
\tightlist
\item
  If \(|w(\gamma)| > 1\): \(|w(\gamma^n)| = |w(\gamma)|^n \to \infty\).
  The amplitude diverges with each iteration.
\item
  If \(|w(\gamma)| < 1\): \(|w(\gamma^n)| = |w(\gamma)|^n \to 0\). The
  amplitude collapses to zero with each iteration.
\end{itemize}

In either case, the probability \(P \propto |w(\gamma^n)|^2\) either
diverges or goes to zero. A diverging probability is unphysical. A
vanishing probability means the entity ceases to exist with probability
1 after finitely many iterations --- inconsistent with a \emph{stable}
persistent loop by assumption.

Therefore \(|w(\gamma)| = 1\) for any loop associated with a stable
entity. Since every bond, particle, and persistent structure in the
cause-plex is modeled as a stable loop {[}Part 1.5: Causors{]}, this
holds for all physically relevant loops. \qed

\begin{corollary}
\label{corollary:3-2}
The image of $w$ lies in the unit sphere of $G$: $w(\gamma) \in G_1 = \{g \in G : |g| = 1\}$ for all closed loops of stable entities.

\end{corollary}

\begin{remark}
\label{remark:3-3}
Proposition 3.1 does not yet constrain the type of complex numbers — it shows only that weights are magnitude-1. The group $G_1$ could still be $\{+1\}$, $\{+1,-1\}$, U(1), SU(2), or more exotic structures. The next two parts narrow this down.

\end{remark}

\begin{remark}
\label{remark:3-4}
The stability condition is physical, not mathematical. In a cause-plex without stable loops — one in which all structures are transient — this argument does not apply. But such a cause-plex supports no persistent entities and therefore no observers. The stable-loop condition is precisely the condition that the cause-plex supports the entity hierarchy [Part 1.5: Causors]. Proposition 3.1 holds in any cause-plex that supports complex auto-causal loops with $\rho_{\mathrm{ac}} > 0$.
\end{remark}

\medskip

\section{Part 2: Causal Continuity Forces a Topological
Group}\label{part-2-causal-continuity-forces-a-topological-group}

\subsection{Single-Event Deformation}\label{single-event-deformation}

The branch graph \(\mathcal{B}\) (Definition 2.4) connects any two
histories that differ by exactly one causal event. A loop \(\gamma\) in
history \(\mathcal{C}_\alpha\) can be ``deformed'' to a nearby loop
\(\gamma'\) in history \(\mathcal{C}_\beta\) if \(\mathcal{C}_\beta\) is
adjacent to \(\mathcal{C}_\alpha\) in \(\mathcal{B}\) and \(\gamma'\)
differs from \(\gamma\) only in the event corresponding to the branch
difference.

\begin{definition}[Single-event deformation]
\label{definition:4-1}
A single-event deformation of loop $\gamma$ is a loop $\gamma'$ obtained by replacing exactly one event $e_k \in \gamma$ with an adjacent event $e_k'$ such that:
1. $e_k'$ and $e_k$ differ by one causal relation in $\mathcal{B}$
2. The resulting sequence $\gamma'$ remains a valid causal loop
\end{definition}

The collection of loops reachable from \(\gamma\) by a finite sequence
of single-event deformations is the \textbf{deformation class}
\([\gamma]\).

\subsection{The Causal Continuity
Condition}\label{the-causal-continuity-condition}

\begin{assumption}[Causal Continuity]
An amplitude assignment $w$ satisfies causal continuity if, for any loop $\gamma$ and any single-event deformation $\gamma'$ of $\gamma$:
$$|w(\gamma') - w(\gamma)| \leq C \cdot d(\gamma, \gamma')$$
for some constant $C$ and a natural distance $d$ on the loop space derived from the branch graph metric.
\end{assumption}

In words: loops that are close in the branch graph have close amplitudes
--- no discontinuous jumps when a single causal event is added or
removed.

\textbf{Status: Postulate.} Condition CC is an independent physical
assumption --- not derivable from the cause-plex primitive, the locality
assumption (L), or the causal dependency axiom (CD). It is the discrete
analogue of the smoothness assumption that in standard QFT is guaranteed
by the manifold structure. It is physically well-motivated (amplitudes
should not jump discontinuously from one causal event), but it is a
genuine additional input to the framework. Theorem 5.2 requires CC
explicitly in its premises, and CC should be understood as a fourth
independent postulate alongside P3, Q', and the stability condition.

In the full postulate inventory for this paper: - \textbf{Stability}
(\(\rho_{\mathrm{ac}} > 0\)): the cause-plex supports stable persistent
loops - \textbf{Locality (L)} and \textbf{CD}: from the spacetime
companion paper - \textbf{Postulate Q'}: the cause-plex is multiway -
\textbf{Postulate CC} (this): amplitude assignments are
Lipschitz-continuous with respect to single-event deformations

\subsection{Causal Simple-Connectedness and Loop Space
Connectivity}\label{causal-simple-connectedness-and-loop-space-connectivity}

The full statement ``all loops in \(\mathcal{C}^{\ast}\) are
deformation-connected'' is false in general --- a cause-plex with a
\emph{causal void} (a region containing no events, bounded by a closed
loop) admits topologically isolated loops whose winding number around
the void is a deformation invariant. No finite sequence of single-event
deformations can change the winding number, so loops with different
winding numbers live in different connected components of the loop
space.

However, this obstruction does not arise in physical cause-plexes
supporting observer-class entities.

\begin{definition}[Causally simply-connected]
\label{definition:4-2}
A cause-plex $\mathcal{C}^{*}$ is causally simply-connected* if every causal region bounded by a closed loop contains at least one causal event — equivalently, there are no causal voids that create non-contractible loops.

\end{definition}

\begin{lemma}[Observer-class cause-plexes are causally simply-connected]
\label{lemma:4-3}
Any cause-plex $\mathcal{C}^{*}$ supporting entities with $\rho_{\mathrm{ac}} > 0$ is causally simply-connected at all scales $\ell \geq \ell_{\mathrm{ac}}$, where $\ell_{\mathrm{ac}}$ is the characteristic scale of the auto-causal loops sustaining those entities. The CI condition ($\mathrm{CI} > \mathrm{CI}_{\min}$) from the Stable Observer Manifold Conjecture strengthens the observer-class requirement but is not needed for this lemma: $\rho_{\mathrm{ac}} > 0$ alone is sufficient.*
\end{lemma}

\begin{proof}

*Step 1: Define the relevant scale.* Let $\ell_{\mathrm{ac}}$ be the smallest spatial scale at which auto-causal loops with $\rho_{\mathrm{ac}} > 0$ operate. Concretely, $\ell_{\mathrm{ac}}$ is the diameter of the smallest recurring causal loop that contributes to $\rho_{\mathrm{ac}}$. For biological entities this is on the order of molecular bond lengths ($\sim 10^{-10}$ m); for fundamental particles, the Compton wavelength ($\sim 10^{-15}$ m or smaller). We assess causal simple-connectedness at scales $\ell \geq \ell_{\mathrm{ac}}$.

*Step 2: $\rho_{\mathrm{ac}} > 0$ implies event density at scale $\ell_{\mathrm{ac}}$.* By definition, $\rho_{\mathrm{ac}}$ measures the density of self-sustaining causal loops per unit volume [Part 1.5: Causors]. A loop of scale $\ell_{\mathrm{ac}}$ returning to its initial state within time $\tau_{\mathrm{ac}}$ requires at least one causal event within each region of diameter $\ell_{\mathrm{ac}}$ per period $\tau_{\mathrm{ac}}$ — otherwise the loop has a segment with no causal events, and causal influence cannot propagate across it (by Definition 1.1: a causal event is the unit of state transition; without events, the state does not evolve, and the loop cannot close). Therefore: in any cause-plex with $\rho_{\mathrm{ac}} > 0$, every bounded causal region of diameter $\geq \ell_{\mathrm{ac}}$ contains at least one causal event.

*Step 3: Event density implies causal simple-connectedness at scale $\ell_{\mathrm{ac}}$.* A causal void that creates a non-contractible loop at scale $\ell$ must be a region of diameter $\geq \ell$ with no causal events (otherwise the loop can be rerouted through the interior events). By Step 2, no such void exists at scale $\ell \geq \ell_{\mathrm{ac}}$. Therefore the loop space at all scales $\ell \geq \ell_{\mathrm{ac}}$ is free of void-induced non-contractible loops. $\mathcal{C}^{\ast}$ is causally simply-connected at the relevant scales. \qed

**Remark 4.3a (Sub-$\ell_{\mathrm{ac}}$ scales).** Below $\ell_{\mathrm{ac}}$, causal voids may exist (e.g., at the Planck scale, local finiteness permits sparse event structure). However, loops at sub-$\ell_{\mathrm{ac}}$ scales are below the resolution of the auto-causal entities in question — no observer-class entity can form or detect amplitude-weight distinctions at that scale. The amplitude assignment $w$ is defined on loops at scales accessible to the entity ($\ell \geq \ell_{\mathrm{ac}}$), so sub-$\ell_{\mathrm{ac}}$ topology does not affect the argument.

**Remark 4.3b (Formal status).** This proof is complete given: (i) the definition of $\rho_{\mathrm{ac}}$ from Part 1.5 (auto-causal loop density — see Definition 2.10 below); (ii) the causal event primitive (Definition 2.1); (iii) the assumption that causal influence requires causal events (no action at a distance without an intervening event). Assumption (iii) is a joint consequence of the locality assumption (L) and the Causal Dependency Axiom (CD) from the spacetime companion paper: L bounds state domains to finite local regions; CD ensures that if $e_2$ reads state written by $e_1$, then $e_1 \prec e_2$ — together these force causal influence to propagate only through chains of events. There are no remaining informal steps.

\end{proof}

\begin{theorem}[Loop space connectivity for observer-class cause-plexes]
\label{theorem:4-4}
If $\mathcal{C}^{*}$ is causally simply-connected, then the loop space $\Omega(\mathcal{C}^{*}, e_{*})$ is connected under single-event deformations: any two closed loops based at $e_{*}$ are connected by a finite sequence of single-event deformations.
\end{theorem}

\textbf{Proof sketch.}

\begin{enumerate}
\def\labelenumi{\arabic{enumi}.}
\item
  \textbf{History space connectivity} (proved in full, see proof
  §2): any two histories
  \(\mathcal{C}_\alpha, \mathcal{C}_\beta \in \mathcal{C}^{\ast}\) are
  connected by a finite sequence of single-relation changes in their
  Hasse diagrams. The branch graph \(\mathcal{B}\) is connected.
\item
  \textbf{Loop tracking along history changes.} Given loops \(\gamma_1\)
  in \(\mathcal{C}_{\alpha_1}\) and \(\gamma_2\) in
  \(\mathcal{C}_{\alpha_2}\), follow the path in \(\mathcal{B}\) from
  \(\mathcal{C}_{\alpha_1}\) to \(\mathcal{C}_{\alpha_2}\). At each step
  --- each single-relation change --- the loop either (a) passes through
  a non-critical event (deforms cleanly, Proposition 4.5 below), or (b)
  passes through a critical event.
\item
  \textbf{Critical events reroute by causal simple-connectedness.} A
  loop-critical event is one with no alternative causal path around it
  within the current history. CSS guarantees every bounded causal region
  contains events; the Critical-Events Rerouting
  proves this is sufficient: via a finite sequence of covering-relation
  additions (each an adjacent-history transition) and one event
  replacement, any loop-critical event can be rerouted while keeping the
  loop closed at every step. The argument uses CSS via Lemma L1
  (intermediate event exists), Lemma L2 (incorporate it via history
  transitions), and Lemma L3 (perform the replacement). The termination
  of sequential rerouting for multiple critical events follows from the
  Noetherian property of locally finite posets (Corollary 5.1 there).
  This step is now fully proved.
\item
  \textbf{Conclusion.} Every step in the path from
  \(\mathcal{C}_{\alpha_1}\) to \(\mathcal{C}_{\alpha_2}\) is
  accompanied by a valid loop deformation. \(\gamma_1\) is deformed
  step-by-step until it arrives in \(\mathcal{C}_{\alpha_2}\), then
  further deformed within \(\mathcal{C}_{\alpha_2}\) (which is
  simply-connected, so all loops there are contractible relative to each
  other) to match \(\gamma_2\). \qed
\end{enumerate}

\begin{proposition}[Non-critical deformations are always valid]
\label{proposition:4-5}
If event $e_k \in \gamma$ is not loop-critical — if at least one alternative causal path connects $e_{k-1}$ to $e_{k+1}$ not through $e_k$ — then any single-event deformation modifying $e_k$ yields a valid closed loop $\gamma'$.
\end{proposition}

\begin{proof}

\end{proof}

\begin{proposition}[Connectivity of $w$-image]
\label{proposition:4-6}
If $\mathcal{C}^{*}$ is causally simply-connected and satisfies Causal Continuity (Condition CC), then the image $w(\Omega(\mathcal{C}^{*}, e_{*}))$ is a connected subset of $G_1$.*
\end{proposition}

\begin{proof}

\end{proof}

\begin{corollary}
\label{corollary:4-7}
Under the conditions of Proposition 4.6, $w(\Omega)$ generates a connected subgroup of $G_1$.
\end{corollary}

\medskip

\section{\texorpdfstring{Part 3: Minimality Forces
\(G = \mathrm{U}(1)\)}{Part 3: Minimality Forces G = \textbackslash mathrm\{U\}(1)}}\label{part-3-minimality-forces-g-mathrmu1}

\subsection{Classification of Connected Compact
Groups}\label{classification-of-connected-compact-groups}

Given Parts 1 and 2, the amplitude weight group \(G_1\) is: - A subgroup
of the unit sphere of some normed division algebra (from Proposition
3.1) - Connected (from Corollary 4.7) - Compact (unit sphere) - Acting
consistently on amplitudes with \(P \geq 0\) (from Definition 2.9)

The normed division algebras are, by Hurwitz's theorem (1898):
\(\mathbb{R}\), \(\mathbb{C}\), \(\mathbb{H}\) (quaternions),
\(\mathbb{O}\) (octonions). Their unit spheres are \(S^0 = \{±1\}\),
\(S^1 = \mathrm{U}(1)\), \(S^3 = \mathrm{Sp}(1) \cong \mathrm{SU}(2)\),
\(S^7\) (not a Lie group).

{\def\LTcaptype{none} % do not increment counter
\begin{longtable}[]{@{}
  >{\raggedright\arraybackslash}p{(\linewidth - 6\tabcolsep) * \real{0.2500}}
  >{\raggedright\arraybackslash}p{(\linewidth - 6\tabcolsep) * \real{0.2500}}
  >{\raggedright\arraybackslash}p{(\linewidth - 6\tabcolsep) * \real{0.2500}}
  >{\raggedright\arraybackslash}p{(\linewidth - 6\tabcolsep) * \real{0.2500}}@{}}
\toprule\noalign{}
\begin{minipage}[b]{\linewidth}\raggedright
Algebra
\end{minipage} & \begin{minipage}[b]{\linewidth}\raggedright
Unit sphere
\end{minipage} & \begin{minipage}[b]{\linewidth}\raggedright
Connected?
\end{minipage} & \begin{minipage}[b]{\linewidth}\raggedright
Gives
\end{minipage} \\
\midrule\noalign{}
\endhead
\bottomrule\noalign{}
\endlastfoot
\(\mathbb{R}\) & \(S^0 = \{+1, -1\}\) & No --- two components & Real QM
--- too weak \\
\(\mathbb{C}\) & \(S^1 = \mathrm{U}(1)\) & \textbf{Yes} & Standard
complex QM \\
\(\mathbb{H}\) & \(S^3 = \mathrm{SU}(2)\) & \textbf{Yes} & Quaternionic
QM --- too strong \\
\(\mathbb{O}\) & \(S^7\) & Yes but not a group & Inconsistent
composition rule \\
\end{longtable}
}

Connectivity (from Part 2) eliminates \(\mathbb{R}\). Octonions are
eliminated because \(S^7\) does not form a group --- the composition
rule (Definition 2.8) requires group structure. This leaves
\(\mathbb{C}\) and \(\mathbb{H}\).

\subsection{Eliminating Real and Quaternionic
QM}\label{eliminating-real-and-quaternionic-qm}

We use two peer-reviewed results to eliminate \(\mathbb{R}\) and
\(\mathbb{H}\) rigorously.

\begin{theorem}\cite{Baez2012}
\label{theorem:5-1}
Quantum theory may be formulated over $\mathbb{R}$, $\mathbb{C}$, or $\mathbb{H}$. Among these, only $\mathbb{C}$ admits both:
1. *Local tomography: the state of a composite system is fully determined by local measurement statistics on subsystems — $\dim(\mathrm{States}(AB)) = \dim(\mathrm{States}(A)) \times \dim(\mathrm{States}(B))$*
2. *The standard tensor product for composite systems: $\mathcal{H}_{AB} = \mathcal{H}_A \otimes \mathcal{H}_B$*

*$\mathbb{R}$-QM and $\mathbb{H}$-QM both fail local tomography. Source: Baez, J.C. "Division algebras and quantum theory." Foundations of Physics 42, 819–855 (2012). DOI: 10.1007/s10701-011-9566-z*

**Experimental confirmation.** Real quantum mechanics ($\mathbb{R}$-QM) makes predictions that differ from complex QM in entanglement experiments involving three or more parties. These differences have been experimentally tested and standard complex QM has been confirmed, ruling out $\mathbb{R}$-QM. Source: Renou et al. "Quantum theory based on real numbers can be experimentally falsified." *Nature* 600, 625–629 (2021). DOI: 10.1038/s41586-021-04160-4

**Application to the cause-plex.** In the cause-plex framework, local tomography is not an independent postulate — it follows from the locality assumption (L): a local observer, with bounded causal influence (finite state domain by Assumption L), can only access local measurement outcomes. If the global state were not determined by local statistics, non-local information would be required to determine the state of a composite system, violating L. Therefore L $\to$ local tomography $\to$ by Theorem 5.1, $G \neq \mathbb{R}$ and $G \neq \mathbb{H}$.
\end{theorem}

The tensor product condition follows from the factorizability of the
cause-plex for non-interacting systems: when
\(D(A) \cap D(B) = \emptyset\) (by Assumption L for spacelike-separated
entities),
\(\mathcal{C}^{\ast}_{AB} = \mathcal{C}^{\ast}_A \times \mathcal{C}^{\ast}_B\),
requiring \(\mathcal{H}_{AB} = \mathcal{H}_A \otimes \mathcal{H}_B\).

\textbf{Status: Proved} (given Baez 2012, peer-reviewed; experimental
confirmation in Renou et al.~2021). This replaces the earlier informal
citation of Aaronson (2004) with rigorous published results.

Therefore quaternions are eliminated, and the unique remaining option is
\(\mathbb{C}\).

\begin{theorem}[Main theorem]
\label{theorem:5-2}
Let $\mathcal{C}^{*}$ be a multiway cause-plex supporting stable persistent loops ($\rho_{\mathrm{ac}} > 0$), satisfying Condition CC (Causal Continuity) and Assumption L (locality). Then $\mathcal{C}^{*}$ is causally simply-connected (Lemma 4.3), and every amplitude assignment $w$ satisfying Definition 2.8 maps closed loops to $\mathrm{U}(1)$:
$$w(\gamma) = e^{i\phi(\gamma)} \quad \text{for some real-valued } \phi: \Omega(\mathcal{C}^{\ast}, e_{\ast}) \to \mathbb{R}$$
\end{theorem}

\begin{proof}
- Proposition 3.1: $|w(\gamma)| = 1$ $\to$ $w(\gamma) \in G_1$ (unit sphere of $G$).
- Lemma 4.3: $\rho_{\mathrm{ac}} > 0$ $\to$ $\mathcal{C}^{\ast}$ is causally simply-connected.
- Theorem 4.4: causal simple-connectedness $\to$ $\Omega(\mathcal{C}^{\ast}, e_{\ast})$ is connected.
- Proposition 4.6 + Corollary 4.7: $w(\Omega)$ is a connected subgroup of $G_1$.
- Hurwitz's theorem: $G \in \{\mathbb{R}, \mathbb{C}, \mathbb{H}, \mathbb{O}\}$.
- $\mathbb{O}$ eliminated: unit sphere $S^7$ is not a group.
- $\mathbb{R}$ eliminated: unit sphere $S^0 = \{+1,-1\}$ is disconnected.
- $\mathbb{H}$ eliminated: Result 5.1 (Aaronson), violated by locality of cause-plex.
- Unique remaining option: $G = \mathbb{C}$, $G_1 = S^1 = \mathrm{U}(1)$.
- Therefore $w(\gamma) = e^{i\phi(\gamma)}$ for real $\phi$. \qed

**Remaining formal step.** The proof of Lemma 4.3 (observer-class $\to$ causally simply-connected) requires formalizing that $\rho_{\mathrm{ac}} > 0$ implies event density above the scale at which non-contractible loops would arise. This is the single remaining gap; see Open Problem 7.1.
\end{proof}

\medskip

\section{\texorpdfstring{From \(\mathrm{U}(1)\) to
\(e^{iS/\hbar}\)}{From \textbackslash mathrm\{U\}(1) to e\^{}\{iS/\textbackslash hbar\}}}\label{from-mathrmu1-to-eishbar}

\subsection{Layer Architecture}\label{layer-architecture}

Theorem 5.2 is a \textbf{Layer 0} result: it derives that amplitude
weights lie in U(1) using only the causal event primitive, the multiway
structure, stability, CSS, and locality. No energy, no action, no
\(\hbar\).

The identification \(\phi(\gamma) = S[\gamma]/\hbar\) is a \textbf{Layer
C} statement: it names what the abstract U(1) phase corresponds to in
the Layer C vocabulary, where time-translation symmetry holds and energy
is defined. This is not a derivation of S from the U(1) result --- it is
a \emph{definitional identification} that locates the Layer 0 structural
fact within the Layer C description. These are different claims;
confusing them was the source of the apparent circularity in the earlier
formulation.

The two claims together form a complete picture:

{\def\LTcaptype{none} % do not increment counter
\begin{longtable}[]{@{}
  >{\raggedright\arraybackslash}p{(\linewidth - 4\tabcolsep) * \real{0.3333}}
  >{\raggedright\arraybackslash}p{(\linewidth - 4\tabcolsep) * \real{0.3333}}
  >{\raggedright\arraybackslash}p{(\linewidth - 4\tabcolsep) * \real{0.3333}}@{}}
\toprule\noalign{}
\begin{minipage}[b]{\linewidth}\raggedright
Layer
\end{minipage} & \begin{minipage}[b]{\linewidth}\raggedright
Claim
\end{minipage} & \begin{minipage}[b]{\linewidth}\raggedright
Status
\end{minipage} \\
\midrule\noalign{}
\endhead
\bottomrule\noalign{}
\endlastfoot
0 & \(w(\gamma) \in \mathrm{U}(1)\) & \(\checkmark\) Proved (Theorem
5.2) \\
C & \(\phi(\gamma) = S[\gamma]/\hbar\) where
\(S = \sum_k \Delta E(e_k) \cdot \tau_{e_k}\) & \(\checkmark\)
Definitional (Proposition 6.2 below) \\
\end{longtable}
}

\subsection{The Phase Function Is
Additive}\label{the-phase-function-is-additive}

\begin{proposition}[Additivity of $\phi$]
\label{proposition:6-1}
The phase function $\phi: \Omega(\mathcal{C}^{*}, e_{*}) \to \mathbb{R}$ satisfies $\phi(\gamma_1 \cdot \gamma_2) = \phi(\gamma_1) + \phi(\gamma_2) \pmod{2\pi}$.
\end{proposition}

\begin{proof}
\end{proof}

An additive real-valued function on paths is an action functional.
Therefore \(\phi\) is an action in the cause-plex sense:
\(\phi(\gamma) = S_\phi[\gamma]\) for some function \(S_\phi\) that
accumulates additively along causal paths.

\subsection{\texorpdfstring{Identifying \(\hbar\) from the
Cause-Plex}{Identifying \textbackslash hbar from the Cause-Plex}}\label{identifying-hbar-from-the-cause-plex}

The quantum of action \(\hbar\) connects the dimensionless U(1) phase to
the dimensionful action \(S\) with units energy \(\times\) time. In the
cause-plex framework, both time and energy are defined relative to a
reference oscillation, and \(\hbar\) emerges from their product at the
minimal scale.

\textbf{Time as normalization.} From Definition 3.2 of the spacetime
companion paper, proper time is a ratio of event counts:
\[\tau(\gamma) = \frac{|\gamma|}{|\mathcal{C}_{\text{ref}}|}\] where
\(|\gamma|\) is the number of causal events along \(\gamma\) and
\(|\mathcal{C}_{\text{ref}}|\) is the number of events per reference
period. Time is dimensionless until the reference period is assigned a
unit. The minimum time step is:
\[\tau_{\min} = \frac{1}{|\mathcal{C}_{\text{ref}}|}\] one event divided
by the number of events per reference period. This is the normalization
constant that converts event counts to physical time --- it is not the
Planck time by definition, but calibrates to the Planck scale when
\(|\mathcal{C}_{\text{ref}}|\) is calibrated accordingly.

\textbf{Energy as Layer C defined quantity.} At Layer C, energy
\(\Delta E\) is the Noether conserved quantity under time-translation
symmetry (spacetime paper §5.2). Its minimum value \(\Delta E_{\min}\)
is the smallest energy difference between states connected by a causal
event --- measured relative to the same reference oscillation.

\textbf{\(\hbar\) as the product.} The phase \(\phi(\gamma)\) is
dimensionless. The action
\(S[\gamma] = \sum_k \Delta E(e_k) \cdot \tau_{e_k}\) has units energy
\(\times\) time. The conversion factor is \(\hbar\):
\[\phi(\gamma) = \frac{S[\gamma]}{\hbar}\]

For the minimum causal event --- one event of energy \(\Delta E_{\min}\)
lasting \(\tau_{\min}\) --- the accumulated action is
\(S_{\min} = \Delta E_{\min} \cdot \tau_{\min}\), and the corresponding
phase is \(\phi_{\min}\) (one minimal phase step). Normalizing so
\(\phi_{\min} = 1\) (one radian per minimal causal event, the natural
unit for U(1)):

\[\hbar = \frac{\Delta E_{\min} \cdot \tau_{\min}}{\phi_{\min}} = \Delta E_{\min} \cdot \tau_{\min} = \frac{\Delta E_{\min}}{|\mathcal{C}_{\text{ref}}|}\]

\textbf{Interpretation.} \(\hbar\) is the ratio of the minimum energy
per causal event to the number of events per reference period. It is not
derived from the cause-plex axioms --- its numerical value is empirical.
What the framework provides is its \emph{structural identity}: \(\hbar\)
is the action quantum set by the discrete minimum-event structure of the
cause-plex, converted to SI units by the same normalization constant
(\(|\mathcal{C}_{\text{ref}}|\)) that defines physical time.

\textbf{Why this is not circular.} The old formulation identified
\(\tau_{\min}\) with the Planck time \(\sqrt{\hbar G/c^5}\), creating
circularity. The present formulation identifies
\(\tau_{\min} = 1/|\mathcal{C}_{\text{ref}}|\) --- a ratio of event
counts, with no reference to \(\hbar\). The connection
\(\hbar = \Delta E_{\min} \cdot \tau_{\min}\) then expresses \(\hbar\)
structurally without presupposing it.

\subsection{The Layer C
Identification}\label{the-layer-c-identification}

\begin{remark}[Phase-action correspondence at Layer C]
\label{remark:6-2}
This is a definitional identification, not a derivation. It assigns Layer C vocabulary (action, energy, $\hbar$) to the abstract U(1) phase established at Layer 0. It names what $\phi$ corresponds to; it does not constrain what $\phi$'s value is for any specific system (electrons, photons, etc.) — that requires a specific Lagrangian, which is not provided here.

*In the Layer C description — where time-translation symmetry holds, energy $E$ is defined as the Noether conserved quantity, and time $\tau$ is defined as event count ratio — the additive phase function $\phi$ of Theorem 5.2 corresponds to:*
$$\phi(\gamma) = \frac{S[\gamma]}{\hbar}, \quad S[\gamma] = \sum_{k=1}^n \Delta E(e_k) \cdot \tau_{e_k}, \quad \hbar = \Delta E_{\min} \cdot \tau_{\min}$$

*In the continuum limit, $S[\gamma] \to \int_\gamma L\, dt$ (the standard Feynman action), and the identification $w(\gamma) = e^{iS[\gamma]/\hbar}$ follows.*
\end{remark}

\begin{proof}

*Step 1: $\phi$ is additive (Proposition 6.1). Any additive real function on causal paths defines an action. The cause-plex action $S[\gamma] = \sum_k \Delta E(e_k) \cdot \tau_{e_k}$ is additive by construction (sums over events). Therefore $\phi(\gamma)$ and $S[\gamma]/\hbar$ are both additive with the same domain and codomain.*

*Step 2: $\hbar$ is the unique scale factor.* Given any two additive real functions on causal paths that agree on the minimal event ($\phi_{\min} = S_{\min}/\hbar$), they are proportional by the same factor at every scale (since both are homomorphisms from the concatenation monoid to $\mathbb{R}$, and the ratio on the generator determines the ratio everywhere). The scale factor is $\hbar = \Delta E_{\min} \cdot \tau_{\min}$ by definition of the minimum action quantum.

*Step 3: Continuum limit.* In the continuum limit ($\tau_{\min} \to 0$, $|\mathcal{C}_{\text{ref}}| \to \infty$ with $\tau_{\min} \cdot |\mathcal{C}_{\text{ref}}| = T_{\text{ref}}$ fixed), the sum $\sum_k \Delta E(e_k) \cdot \tau_{e_k}$ becomes $\int L\, dt$ by the standard Riemann-sum argument, where $L = T - V$ is the classical Lagrangian. This is not a new result — it is the standard derivation of the path integral from the discrete approximation \cite{Feynman1965}. \qed

**Layer note.** Proposition 6.2 is a Layer C result. The energy $\Delta E$ and the Lagrangian $L$ are Layer C quantities — defined where time-translation symmetry holds. The identification has no content at Layer 0 because energy is not defined there. This is not a limitation; it correctly reflects that $e^{iS/\hbar}$ is a Layer C description of the Layer 0 fact $w(\gamma) \in \mathrm{U}(1)$.

**On the AUDIT.md Priority 1a circularity.** The earlier formulation claimed to derive energy from Noether applied to the cause-plex action, while defining the action using energy. This was circular. The present formulation dissolves the circularity: the Layer 0 fact (U(1) amplitudes) is proved without energy; the Layer C vocabulary (S, E, $\hbar$) is assigned to name that fact at the appropriate layer. No derivation of energy from the action is claimed or needed.
\end{proof}

\medskip

\section{Open Problems}\label{open-problems}

\subsection{\texorpdfstring{\st{Open Problem 7.1} ---
Resolved}{Open Problem 7.1 --- Resolved}}\label{open-problem-7.1-resolved}

Lemma 4.3 is now formally proved (Section 4.3). The proof proceeds in
three steps: (1) define \(\ell_{\mathrm{ac}}\) as the characteristic
auto-causal loop scale; (2) show \(\rho_{\mathrm{ac}} > 0\) implies
event density at scale \(\ell_{\mathrm{ac}}\) (otherwise auto-causal
loops cannot close); (3) show event density implies causal
simple-connectedness at that scale (no void large enough to create
non-contractible loops exists).

The proof relies on the locality assumption (L) from the spacetime
companion paper, the definition of \(\rho_{\mathrm{ac}}\) from Part 1.5,
and the causal event primitive. No additional postulates are required.

\textbf{For completeness} (see proof
- The full loop-space connectivity conjecture is false in general ---
causal voids create isolated loops with invariant winding number. - For
causally simply-connected cause-plexes (satisfied by all observer-class
cause-plexes by Lemma 4.3), Theorem 4.4 holds. - History space
connectivity is fully proved (proof attempt §2).

\textbf{Status: Lemma 4.3 resolved.} The Critical-Events Rerouting Lemma
(Theorem 4.4 Step 3, formerly Open Problem 7.3) is also proved --- see
Theorem 5.2 is now complete subject only to Conjecture 6.2.

\medskip

\subsection{\texorpdfstring{\st{Open Problem 7.2} --- Resolved
(Proposition
6.2)}{Open Problem 7.2 --- Resolved (Proposition 6.2)}}\label{open-problem-7.2-resolved-proposition-6.2}

The phase-action identification is proved as Proposition 6.2 (Section
6.4). The key reframing: this is a Layer C definitional identification,
not a Layer 0 derivation. The earlier circularity (energy used to define
action, action used to derive energy) is dissolved by correctly locating
each claim at its layer.

Three steps in the proof: (1) additivity of \(\phi\) (from Definition
2.8); (2) \(\hbar\) as the unique scale factor between the dimensionless
phase and the dimensionful action, expressed non-circularly as
\(\hbar = \Delta E_{\min}/|\mathcal{C}_{\text{ref}}|\); (3) the
continuum limit \(\sum \Delta E \cdot \tau \to \int L\,dt\) by standard
Riemann-sum argument.

\textbf{Status: Closed.} The full result
\(w(\gamma) = e^{iS[\gamma]/\hbar}\) is now established:
\(w(\gamma) \in \mathrm{U}(1)\) at Layer 0 (Theorem 5.2), and
\(\phi(\gamma) = S[\gamma]/\hbar\) at Layer C (Proposition 6.2).

\medskip

\subsection{\texorpdfstring{\st{Open Problem 7.3} ---
Resolved}{Open Problem 7.3 --- Resolved}}\label{open-problem-7.3-resolved}

The Critical-Events Rerouting Lemma is proved in the companion document
The proof has three sub-lemmas:

\begin{itemize}
\tightlist
\item
  \textbf{L1:} CSS guarantees an intermediate event \(c\) exists between
  the endpoints of any loop-critical covering relation, in some history
  of \(\mathcal{C}^{\ast}\).
\item
  \textbf{L2:} \(c\) can be incorporated into the ambient history via
  two covering-relation additions (each a valid adjacent-history
  transition), preserving loop validity throughout.
\item
  \textbf{L3:} In the enriched history, the loop-critical event can be
  replaced by \(c\) via a single-event deformation (Case 1: direct
  connection; Case 2: transitive via an additional history step).
\end{itemize}

Sequential rerouting of multiple critical events terminates by the
Noetherian property of locally finite posets (Corollary 5.1).

\textbf{Status: Substantially argued.} The step3\_lemma closes the main
gap but retains three sub-gaps (Lemma L1 region assumption, Assumption
R, Case 2 global consistency). These are closeable; the approach is
directionally correct. Theorem 5.2 follows given these sub-gaps.
Identification 6.2 establishes the Layer C naming of φ as S/\(\hbar\).

\medskip

\section{Discussion}\label{discussion}

\subsection{What the paper
establishes}\label{what-the-paper-establishes}

Taking stock of what is and is not proved:

{\def\LTcaptype{none} % do not increment counter
\begin{longtable}[]{@{}
  >{\raggedright\arraybackslash}p{(\linewidth - 2\tabcolsep) * \real{0.5000}}
  >{\raggedright\arraybackslash}p{(\linewidth - 2\tabcolsep) * \real{0.5000}}@{}}
\toprule\noalign{}
\begin{minipage}[b]{\linewidth}\raggedright
Claim
\end{minipage} & \begin{minipage}[b]{\linewidth}\raggedright
Status
\end{minipage} \\
\midrule\noalign{}
\endhead
\bottomrule\noalign{}
\endlastfoot
Stable loop amplitudes have unit magnitude & \(\checkmark\) Proved
(Proposition 3.1) \\
Causal continuity implies topological group structure & \(\checkmark\)
Proved given Condition CC \\
History space connected (any two histories reachable) & \(\checkmark\)
Proved (proof attempt §2) \\
Causal voids create isolated loops (counterexample) & \(\checkmark\)
Established (causal annulus) \\
Observer-class cause-plexes are causally simply-connected &
\(\checkmark\) Proved (Lemma 4.3) \\
Loop space connected for CSS cause-plexes & \(\checkmark\) Proved
(Theorem 4.4 + Step 3 Lemma; all sub-gaps closed) \\
Connected compact group over normed division algebra & \(\checkmark\)
Follows via Hurwitz \\
\(\mathbb{R}\) eliminated & \(\checkmark\) Proved --- disconnected unit
sphere \(S^0\) \\
\(\mathbb{H}\) eliminated & \(\checkmark\) Proved --- Baez
\cite{Baez2012} + Renou et al.~(2021, experimental) \\
\(G = \mathrm{U}(1)\) uniquely & \(\checkmark\) Proved (Theorem 5.2) \\
\(\phi = S/\hbar\) at Layer C & \(\checkmark\) Definitional
identification (Identification 6.2) --- names the phase; value requires
specific Lagrangian \\
\(\hbar = \Delta E_{\min}/\lvert\mathcal{C}_{\text{ref}}\rvert\)
non-circular & \(\checkmark\) Derived from event-count time
definition \\
Full result \(w(\gamma) = e^{iS/\hbar}\) & \(\checkmark\) Proved (Layer
0 structural + Layer C naming) \\
\end{longtable}
}

\subsection{Why this matters}\label{why-this-matters}

Given Lemma 4.3 (formally proved), the imaginary unit \(i\) in quantum
mechanics is no longer postulated --- it is derived from:

\begin{enumerate}
\def\labelenumi{\arabic{enumi}.}
\tightlist
\item
  The stability of persistent loops (physical)
\item
  The connectedness of the loop space (mathematical, about the structure
  of \(\mathcal{C}^{\ast}\))
\item
  The locality of measurements (physical)
\end{enumerate}

These three conditions are all prior to quantum mechanics. They are
properties of the cause-plex structure that would be required for a
physically sensible theory regardless of whether quantum mechanics is
true. The derivation would therefore be genuinely foundational --- not
``given QM, here is why \(i\) appears'' but ``given these structural
requirements on any causal theory, QM with complex amplitudes is
forced.''

This is a different and stronger result than Sorkin's quantum measure
theory argument (which derives \(i\) from the structure of probability
theory, not from physical causal structure) or the Wick rotation
argument (which derives \(i\) by analogy with signature mismatch, not
from group theory).

\subsection{Connection to Sorkin}\label{connection-to-sorkin}

Sorkin's grade-2 quantum measure argument and the loop-phase argument
are complementary, not competing. Sorkin asks: given a multiway history
structure, what is the minimal consistent probability theory? The answer
is grade-2, which forces complex-valued decoherence functionals. The
present paper asks: given stable loops in a cause-plex, what amplitude
group is forced? The answer is U(1). These are the same conclusion
reached by different routes --- probability-theoretic and structural
respectively --- providing independent convergent support.

\subsection{Connection to the entity
taxonomy}\label{connection-to-the-entity-taxonomy}

The loop-phase argument is natural within the epimechanics framework
because persistent entities \emph{are} loops. The connection between
quantum mechanics and the entity taxonomy is therefore tight: the same
loop structure that defines bonds, particles, and auto-causal entities
(Part 1.5) is the structure that forces the complex phase weights.
Quantum mechanics is not an additional layer imposed on the cause-plex
--- it is the amplitude theory forced by the loop structure native to
persistent entities.

\medskip

\section{Scope and Limitations}\label{scope-and-limitations}

\subsection{What This Paper
Establishes}\label{what-this-paper-establishes}

This paper is a \textbf{necessary-conditions structural result}. It
shows that \emph{if} a multiway cause-plex satisfies four postulates
(stability, locality/CD, Q', and CC), \emph{then} amplitude weights must
lie in U(1). This is a constraint on the form of quantum mechanics ---
complex phases are forced, not chosen --- but it is not a derivation of
any specific physical prediction.

Concretely, the paper establishes:

{\def\LTcaptype{none} % do not increment counter
\begin{longtable}[]{@{}
  >{\raggedright\arraybackslash}p{(\linewidth - 2\tabcolsep) * \real{0.5000}}
  >{\raggedright\arraybackslash}p{(\linewidth - 2\tabcolsep) * \real{0.5000}}@{}}
\toprule\noalign{}
\begin{minipage}[b]{\linewidth}\raggedright
What is established
\end{minipage} & \begin{minipage}[b]{\linewidth}\raggedright
What is not
\end{minipage} \\
\midrule\noalign{}
\endhead
\bottomrule\noalign{}
\endlastfoot
Amplitude weights \(w(\gamma) \in \mathrm{U}(1)\) & What \(\phi\) is for
electrons, photons, quarks \\
The imaginary unit is structurally forced & Any specific Hamiltonian or
Lagrangian \\
U(1) is the unique minimal continuous group & The Hilbert space for any
physical system \\
\(\phi = S/\hbar\) names the phase at Layer C & Why QED has the specific
coupling constant \(e\) \\
\(\hbar\) is non-circularly identified & Any numerical prediction \\
\end{longtable}
}

\subsection{What Is Needed for Computable
Predictions}\label{what-is-needed-for-computable-predictions}

To compute a physical number from the cause-plex framework requires,
beyond this paper:

\begin{enumerate}
\def\labelenumi{\arabic{enumi}.}
\item
  \textbf{Specific Lagrangians} for each field (QED, QCD, gravity).
  These encode particle content, masses, and couplings. The cause-plex
  framework does not yet derive them --- it establishes only that
  \emph{some} additive phase function exists.
\item
  \textbf{Field operators and Hilbert space structure.} The bridge from
  discrete causal events to quantum fields (\(\hat{\psi}(x)\),
  \(\hat{A}_\mu(x)\)) has not been made. Second quantization applied to
  specific classical fields would need to emerge from the
  coarse-graining of the cause-plex.
\item
  \textbf{The correct path integral measure} (not just the weight type).
  Even given \(w(\gamma) = e^{iS/\hbar}\), the measure
  \(\mathcal{D}\gamma\) over path space is a separate object with its
  own subtleties (Wiener measure, lattice regularization, etc.).
\item
  \textbf{Renormalization.} The cause-plex has a natural UV cutoff at
  \(\ell_{\mathrm{ac}}\), but the explicit renormalization group flow
  connecting Planck-scale discreteness to low-energy observables is not
  worked out.
\end{enumerate}

\subsection{Falsifiability}\label{falsifiability}

In its current form, this paper makes no falsifiable prediction
distinguishable from standard quantum mechanics. The result is
conditional: \emph{if} the postulates hold, \emph{then} U(1). Standard
QM experiments confirm U(1) as expected; they do not test the
postulates.

\textbf{Paths to falsifiable predictions (unworked):}

\begin{itemize}
\item
  \textbf{Planck-scale dispersion.} If \(\ell_{\mathrm{ac}}\) is finite
  and specifiable for fundamental particles, corrections to the
  dispersion relation \(E^2 = p^2c^2 + m^2c^4\) of order
  \(\ell_{\mathrm{ac}} E^3/\hbar c\) are expected. Gamma-ray burst
  timing (Fermi-LAT) constrains such corrections. Requires specifying
  \(\ell_{\mathrm{ac}}\) for photons --- not done here.
\item
  \textbf{CSS violation signatures.} If the CSS condition fails at some
  scale, amplitude weights would acquire additional topological quantum
  numbers (winding sectors). This would predict discrete anomalies in
  interference experiments at that scale. Completely untestable with
  current technology.
\item
  \textbf{Decoherence structure.} The branching structure of
  \(\mathcal{C}^{\ast}\) may impose constraints on decoherence rates or
  preferred basis selection beyond what standard QM predicts. Not
  currently worked out.
\end{itemize}

A companion phenomenological paper specifying the cause-plex at particle
scales would be needed to generate testable predictions.

\subsection{Relationship to Standard
QFT}\label{relationship-to-standard-qft}

This paper sits \emph{below} standard QFT in the foundational hierarchy.
QFT takes complex amplitudes and Hilbert space structure as given; this
paper argues for why those structures are forced. The gap between this
paper's conclusion and any QFT calculation is large --- it is, roughly,
the entire edifice of quantum field theory.

This does not diminish the result: a genuine derivation of why complex
amplitudes must exist (rather than postulating them) is foundational
progress. But it should not be mistaken for a derivation of QFT itself.

\medskip

\section{Conclusion}\label{conclusion}

We have shown that the imaginary unit in quantum mechanics is
structurally forced --- not postulated --- in any multiway cause-plex
satisfying four conditions: stability of persistent loops, causal
continuity (Postulate CC), locality (L + CD), and multiway structure
(Q'). The proof chain is:

\begin{figure}[htbp]
\centering
\fbox{\parbox{0.85\textwidth}{\centering\vspace{2cm}\textit{[Figure 1: Diagram placeholder --- export Mermaid to PDF]}\vspace{2cm}}}
\caption{Figure 1: See Markdown source for diagram definition.}
\label{fig:causeplex_loop_phase-1}
\end{figure}

\begin{enumerate}
\def\labelenumi{\arabic{enumi}.}
\tightlist
\item
  Stability \(\to\) \(|w| = 1\) (Proposition 3.1 --- \(\checkmark\)
  proved)
\item
  Observer-existence \(\to\) causal simple-connectedness (Lemma 4.3 ---
  \(\checkmark\) proved)
\item
  CSS + causal continuity \(\to\) loop space connected \(\to\) \(G\) is
  a connected compact group (Theorem 4.4 --- \(\checkmark\) proved; see
\item
  Hurwitz \(\to\) \(G \in \{\mathbb{C}, \mathbb{H}\}\) (\(\checkmark\)
  proved)
\item
  Locality \(\to\) \(G \neq \mathbb{R}, \mathbb{H}\) (\(\checkmark\)
  proved via Baez 2012, Renou et al.~2021)
\item
  Unique conclusion: \(G = \mathrm{U}(1)\),
  \(w(\gamma) = e^{i\phi(\gamma)}\) (Theorem 5.2 --- \(\checkmark\)
  proved)
\end{enumerate}

All open problems are resolved. The full result
\(w(\gamma) = e^{iS[\gamma]/\hbar}\) is established across two layers:
the structural U(1) fact at Layer 0 (Theorem 5.2), and the physical
action identification at Layer C (Proposition 6.2). The energy-action
circularity of earlier formulations is dissolved by the layer
separation: energy is not derived from the action at Layer 0; rather,
the Layer 0 U(1) phase is \emph{named} as \(S/\hbar\) at Layer C where
energy is well-defined. \(\hbar\) is identified as
\(\Delta E_{\min}/|\mathcal{C}_{\text{ref}}|\) --- the minimum energy
per causal event divided by the reference event count --- with no
circularity.

The full counterexample analysis, proof of history-space connectivity,
and detailed proof sketch for Theorem 4.4 are in the companion document

\medskip

\bibliographystyle{plain}
\bibliography{causeplex_loop_phase}

\end{document}
